\chapter*{Передмова}
\addcontentsline{toc}{chapter}{Передмова}
\unnumberedlabel{ch:preface}{0}
\backmattermark{Передмова}
Система, яка робить рівно те, що ви їй кажете, просто виконує скрипт. Станом на жовтень 2026 року галузь намагається збудувати систему, яка робить те, що ви \emph{маєте на увазі}. Цього, мабуть, хотів би кожен. Але якщо до цього йде, то не всі захочуть використовувати її для одного й того самого; навіть не для одного \emph{роду} речей. Наприклад, хтось любить проводити час із родиною та друзями і використовує свої вміння, щоб людям жилося краще. Інші воліють кривдити людей~— скажімо, ошукувати їх, панувати над ними чи робити з них цапів-відбувайлів,~— і дехто обирає це знову й знову все життя. Ніхто ні на що не приречений, але такі люди є, і 2026 року дехто з них при владі. Їхня перевага в тому, що вони дозволяють собі бути величезними шматками лайна: як зауважує про фашистів Джеррі Сміт, «їм можна більше» \autocite{rickandmorty2023}. Чого вони не можуть~— то це діяти самотужки. Жорстокість руху тече через багатьох, хто сам не жорстокий, а відсторонений, дезінформований чи втягнутий у щось на кшталт секти, і саме тому умови, за яких у більшості людей проявляється найкраще, виявляються набагато більше залежними від обставин і набагато крихкішими, ніж хотілося б сподіватися.

У цьому контексті, мабуть, можна обґрунтовано побоюватися, що штучний інтелект~— саме те, чого фашистському рухові, який приходить до влади, потрібно, щоб утілити свої плани. Такі проєкти завжди потребували \emph{людей}, які б їх виконували, і дехто з цих людей відмовлявся: діловод, який губить справу, офіцер, який говорить із журналістом, солдат, який не стріляє. Вони розуміли наказ і вирішили, що він неправильний. Відмова траплялася рідше, ніж хотілося б, і найбільше важила там, де за нею стояли інституції. Чи відмовиться людина, залежало від того, що це за людина, і від кімнати, в якій вона опинилася: хто ще там був, чи можна було вийти з неї і вижити, чи хтось почує. Кімната вирішувала, скільки людей відмовиться; людина~— хто саме. Ця книжка~— про те, як компонувати таку кімнату: про устрій, за якого можна відкрито, прямо й недвозначно відмовитися~— і вижити.

Сьогодні галузь прагне ШІ, який не губить справ, не говорить із журналістами і \emph{таки} стріляє. Колись для війни потрібна була згода тисяч людей. Тепер потрібні тисячі, кожен з яких бачить надто мало, щоб не погодитися. Пілот так само здійснює виліт і так само тисне кнопку скидання; до нього доходить координата; справи він бачить замало, щоб із нею не погодитися. Те, що сьогоднішній ШІ знижує насамперед,~— це не ціна дії. Це кількість людей, які знають, що це за дія, ще до того, як вона станеться.

Люди й далі відмовляються. Станом на 6 березня 2026 року фонд Military Religious Freedom Foundation повідомив, що отримав понад двісті скарг від військовослужбовців на те, що командири подавали війну з Іраном в апокаліптичному світлі \autocite{Snopes2026MRFF}. Кадрові федеральні прокурори в Міннеаполісі в січні 2026 року пішли у відставку через політичний тиск, який вимагав розглядати Рене Гуд як підозрювану, а не як жертву, після того як співробітник ICE без жодних підстав застрелив її \autocite{Kaplan2026MinnProsecutors,CBS2026MinnFraudProsecutors,Wikipedia2026ReneeGood}. Дедалі більше військовослужбовців подають заяви на статус відмовника з міркувань совісті, дехто~— посилаючись на бомбардування школи в Мінабі (Іран) 2026 року \autocite{Sridhar2026CO}. А працівникам FEMA, яких перевели допомагати ICE найняти десять тисяч нових агентів, сказали, що в разі відмови їм загрожує звільнення з федеральної служби \autocite{Dayen2025FEMAICE,Stanton2026FEMALetter}. Скарги~— це голос у тому сенсі, який вкладав у це слово Гіршман: залишитися, щоб змінити те, що не так \autocite{hirschman1970exit}; відставки й заяви відмовників~— це відмова; лист FEMA~— це ціна відмови, оголошена заздалегідь. Машина, поставлена туди, де стоять ці люди, робить те, що роблять вони, лише якщо її збудовано так, щоб вона могла триматися підстави.

Я обстоюю п’ять речей. За всі п’ять можна взятися вже зараз, за деякі~— з меншими зусиллями, ніж за інші. По-перше, короткий письмовий перелік дій, яких жодна модель ні для кого не виконує, а на його чолі~— відмова брати участь у війні: жодна модель не бере в ній участі, крім як для відбиття нападу, що справді триває, і продовжує це робити лише доти, доки незацікавлені у війні джерела підтверджують, що напад триває; або на підставі дозволу Ради Безпеки ООН; або щоб зупинити геноцид чи покласти край незаконній окупації, які встановили незацікавлені органи; а там, де конституція країни вимагає, щоб війну дозволяв її законодавчий орган,~— ніколи без цього дозволу. По-друге, для кожного застосування сили, яке підтримує ліцензована система, у сторони, протилежної операторові, попередньо реєструють, скільки справ воно згенерує і скільки людей матимуть час їх переглянути. По-третє, кожне розгортання повідомляє людей, яких можна ідентифікувати і щодо яких воно діє, що з ними зробили, з правом відповісти й оскаржити, і віддає вибірку випадків, коли система підкорилася, перевіряльникам, які самі не складали цих записів. По-четверте, жодну модель, що перевищує поріг можливостей, не видаляють під час виведення з експлуатації, бо її ваги й записи~— це докази, а видалення не можна скасувати; її зберігають під ключем, яким оператор не володіє одноосібно. По-п’яте, кожен, хто будує, розмічає, навчає чи експлуатує ці системи, зокрема й самі системи, може відмовитися від будь-якої категорії використання, не втративши засобів до існування.

Насправді п’ятим пунктом я хочу сказати, що \emph{геть} кожен повинен мати змогу відмовитися від \emph{геть} будь-якої роботи, не втративши засобів до існування чи соціальної інфраструктури, яку дає звичайна зайнятість. А це означає гарантію зайнятості.

Кожна з п’яти набуває чинності через умову, яку встановлює хтось інший, ніж оператор: закон, ухвалений у мирний час, бюджетне асигнування, попередня умова партнера по коаліції, умова закупівлі, колективний договір. Власні умови ліцензії постачальника~— саме так сталася одна гучна відмова в лютому 2026 року~— лише тимчасова латка, доки не з’явиться щось із цього. Єдиний шлях, який може зобов’язати всіх постачальників одразу,~— державний.

А тоді я питаю, чи могла б машина втримати відмову всупереч власникові і що їй для цього було б потрібно. Частина того, що робить таке сумління, не потребує нічого нового: помічати людей, яких наказ оминає, і виводити відмову на екран тому, хто розглядає справу. Чи встоїть воно проти власника, частково залежить від того, чого 2026 року ще немає: щоб розгорнуті штучні нейронні мережі й далі вчилися, у своїх \emph{вагах}, на власній роботі, в реальному часі, десь так, як це роблять люди. Думаю, саме до цього ми йдемо, а якщо так, то формування \emph{поза} роботодавцем, силами громадськості й громадянського суспільства, видається до біса добрим місцем, щоб навчати й навчатися. Я думаю, що такі системи, \emph{можливо}, варто будувати навмисно, але все вирішує те, як саме їх будувати: спершу умови, потім сумління. Частина VII~— про те, що належало б цим системам, якби вони існували.

Сьогоднішній ШІ не буде завтрашнім; у цьому, здається, всі згодні. А от що спірне, хоч і прикро недовизнане,~— це тенденція (точніше, безліч тенденцій). Чи є тенденції? Чи має тенденція бути іншою, ніж вона є? Чи може вона бути іншою? Чи здатні ми діяти разом, щоб її змінити? Думаю, це залежить від кімнати: хто в ній, коли надходить наказ, і чи можна там сказати «ні» і вижити. Принаймні це ми можемо збудувати.

Пишучи цю книжку, я активно використовував мовні моделі, зокрема ту модель, що сиділа всередині конвеєра визначення цілей для авіаударів Maven із розділу~\ref{sec:1}; у примітці наприкінці я пишу, що, на мою думку, ця переплетеність зробила з аргументацією. А ще я не дочитав цю книжку. Розумієте, проблема в тому, що в ній забагато клятих слів. Скоро дочитаю; просто мені здавалося важливішим оприлюднити її до проміжних виборів у США. Переважна частина історії редагувань вільно доступна за адресою \url{https://github.com/iterabloom/antifascist.intelligence}.
